
We opened with a puzzle: an agent passes 95\% of tests---but did it make 10 strategic root-cause fixes or 100 random modifications? Existing benchmarks can't tell. Our work in controlled settings shows that strategic debugging can be measured via fix-impact-ratio, a metric discriminating agents that target root causes (ratio 3.90: one fix resolves 3-4 test failures) from those using sequential trial-and-error (ratio 1.07: one fix per failure). Experiments reveal strategic debugging operates as a two-stage feedback loop---agents cluster errors by shared characteristics at $2\times$ random rate (coefficient 1.909, p=0.001), then prioritize fixes yielding 2.$15\times$ higher proportion of high-impact modifications (42.5\% vs 19.8\%)---but transfer to unseen test cases fails (slope ratio 0.82 < 1.5, p=0.504). This clarifies that strategic debugging is a feedback loop effective within revealed test execution, not a learning system enabling autonomous prediction without test execution.

\subsubsection{Summary}

We addressed the gap between outcome-focused metrics (pass@k) and process-based evaluation by introducing fix-impact-ratio, a framework measuring how efficiently agents debug through error feedback. Our main contributions are:

\begin{enumerate}
\item \textbf{Fix-impact-ratio framework with three execution-based metrics} --- fix-impact-ratio (tests per modification), clustering coefficient (pattern recognition), held-out slope (generalization attempt) --- validated with large effect size (Cohen's d=3.07, p=0.0001) demonstrating strong discriminative power on controlled data.
\end{enumerate}

\begin{enumerate}
\item \textbf{Experimental validation of two-stage mechanism} --- agents cluster errors by type (coefficient 1.909 vs random 0.950, p=0.001), enabling prioritization of high-impact fixes (42.5\% vs baseline 19.8\%, 2.$15\times$ improvement), explaining how strategic debugging achieves ratio 3.90 vs sequential 1.07. Transfer learning failed (slope ratio 0.82, p=0.504), refuting pattern-based generalization.
\end{enumerate}

\begin{enumerate}
\item \textbf{Theoretical clarification} --- strategic debugging is a feedback loop (error $\rightarrow$ cluster $\rightarrow$ prioritize $\rightarrow$ fix) effective when error messages guide clustering, not a learning system (observe $\rightarrow$ extract $\rightarrow$ predict) supporting transfer to unseen tests without execution. Framework measures iterative debugging with test access, complementing pass@k correctness metrics.
\end{enumerate}

\subsubsection{Future Directions}

\textbf{Validating Ecological Validity:} Mock implementation established metric sensitivity---fix-impact-ratio can detect strategic behavior when engineered (d=3.07). Whether production agents exhibit clustering at coefficient > 0.3 rates on real competitive programming problems remains unverified. Immediate extension: replicate experiments with production agents and curated datasets to confirm mock results generalize.

\textbf{Understanding Transfer Failure:} Pattern memory extracted patterns from revealed test failures but achieved 0\% usage rate on held-out tests. Two competing explanations: (1) pattern quality insufficient, or (2) transfer fundamentally information-limited. Test with oracle error patterns---if held-out slope ratio > 1.5 with oracle patterns, transfer limitation is pattern extraction quality; if ratio remains < 1.5 even with perfect patterns, predictive debugging without execution is fundamentally constrained.

\textbf{Extending to Real-World Scenarios:} Framework applies to multi-test scenarios (competitive programming, issue solving when test suites exist). Extension: analyze patches with test suites, compute fix-impact-ratio and clustering coefficient, compare agent strategies on production codebases vs synthetic competitive programming.

\textbf{Alternative Mechanism Variants:} Current work tested error-type clustering; alternative clustering dimensions worth exploring: semantic similarity of error messages, code location, computational failure mode. Each may reveal additional mechanism variants beyond error-type clustering.

\textbf{Alternative Transfer Mechanisms:} Current pattern memory failed. Alternative designs worth exploring: few-shot learning, meta-learning across problems, causal code-behavior modeling. Each requires different agent capability than post-hoc clustering.

As code generation agents evolve from single-attempt generators to iterative problem-solvers, understanding how they improve becomes as critical as whether they succeed. Fix-impact-ratio establishes a measurement approach ready for real-world validation, proposing benchmark design beyond pass@k paradigm to measure debugging efficiency---a step toward comprehensive assessment of agentic capability in autonomous software development.
